%=====================================================================
\chapter{Hardware Support for Virtualization}\label{ch:hardware}
%=====================================================================
\locator{2}{Part II --- Hypervisors and Compute Virtualization}

\begin{outcomes}
By the end of this chapter you will be able to:
\begin{tight}
  \item \textbf{Explain} how VMX root and non-root operation make x86 satisfy
        the Popek--Goldberg condition without changing a single instruction's
        encoding.
  \item \textbf{Describe} what the VMCS holds and \textbf{account} for the cost
        of a VM exit in terms of the work the processor must do.
  \item \textbf{Count} the memory accesses in a nested page walk, and
        \textbf{explain} why two-dimensional paging replaced shadow page tables
        despite being slower on a TLB miss.
  \item \textbf{State} what an IOMMU does and why device passthrough is unsafe
        without one.
  \item \textbf{Diagnose} a performance problem as an exit-rate problem from
        measured counters, and \textbf{name} the feature that addresses it.
\end{tight}
\end{outcomes}

\begin{prereq}
\chref{vmm} above all --- the Popek--Goldberg condition, the seventeen
offending instructions, and why trap-and-emulate could not be used. This
chapter is the hardware answer to that chapter's problem, and almost nothing in
it makes sense without it. You also need virtual memory and page tables from
Operating Systems.
\end{prereq}

\begin{keyterms}
VT-x $\bullet$ AMD-V $\bullet$ SVM $\bullet$ VMX root operation $\bullet$ VMX
non-root operation $\bullet$ VMCS $\bullet$ VMCB $\bullet$ VM entry $\bullet$
VM exit $\bullet$ exit reason $\bullet$ world switch $\bullet$ shadow page
table $\bullet$ extended page tables $\bullet$ nested page tables $\bullet$
two-dimensional paging $\bullet$ nested page walk $\bullet$ VPID $\bullet$
IOMMU $\bullet$ VT-d $\bullet$ DMA remapping $\bullet$ interrupt remapping
\end{keyterms}

\section{What Intel and AMD Had to Fix}

\chref{vmm} left the architecture in an awkward position. Seventeen
instructions were sensitive but not privileged, so trap-and-emulate was
unsound; the two workarounds --- binary translation and paravirtualization ---
each cost something substantial. Intel shipped VT-x in 2005 and AMD shipped
AMD-V in 2006, and both took the same approach.

The obvious fix would have been to make the seventeen instructions privileged.
Both vendors rejected it, for a reason worth understanding: it would break
existing software. A user program that executes \texttt{POPF} today and has its
write to \texttt{IF} quietly ignored would begin faulting, and backward
compatibility is the entire commercial basis of the x86 architecture.

So instead of changing what the instructions do, they changed \emph{the mode
the processor can be in}.

\begin{definitionbox}[VMX Root and Non-Root Operation]
VT-x adds a new dimension to the privilege model. Orthogonally to rings 0--3,
the processor is in one of two \term{operations}:

\begin{tight}
  \item \term{VMX root operation} --- what the hypervisor runs in. Behaves like
        an ordinary processor, plus a set of new instructions
        (\texttt{VMXON}, \texttt{VMLAUNCH}, \texttt{VMRESUME}, \texttt{VMREAD},
        \texttt{VMWRITE}) for managing guests.
  \item \term{VMX non-root operation} --- what a guest runs in. All four rings
        are available, so a guest kernel runs in \textbf{ring~0} exactly as it
        expects. But certain instructions and events cause a \term{VM exit}:
        control transfers to the hypervisor in root operation.
\end{tight}

\smallskip
The guest kernel is no longer deprivileged. It runs in ring~0, reads its own
privilege level as 0, and sees the descriptor tables it installed. Ring
aliasing is gone, the seventeen instructions trap, and equivalence holds ---
without a single instruction changing its meaning.
\end{definitionbox}

\dsfig{%
\begin{tikzpicture}[font=\scriptsize,
  mode/.style={rounded corners=3pt, line width=1.1pt, align=center,
               minimum width=5.6cm, minimum height=2.5cm},
  b/.style={rounded corners=2pt, line width=0.7pt, align=center,
            minimum width=4.6cm, minimum height=0.52cm, inner sep=2pt,
            font=\scriptsize},
  lb/.style={font=\scriptsize\itshape, align=center}
]
  \node[mode, draw=primary, fill=primary!6] (R) at (0,0) {};
  \node[font=\small\bfseries, text=primary] at (0,1.0) {VMX root operation};
  \node[b, draw=primary, fill=white, text=primary] at (0,0.35) {hypervisor (ring 0)};
  \node[b, draw=primary!50, fill=white, text=primary!70] at (0,-0.3)
       {\texttt{VMLAUNCH} \ \texttt{VMRESUME}};
  \node[b, draw=primary!50, fill=white, text=primary!70] at (0,-0.92)
       {\texttt{VMREAD} \ \texttt{VMWRITE}};

  \node[mode, draw=secondary, fill=secondary!6] (N) at (9.4,0) {};
  \node[font=\small\bfseries, text=secondary!80!black] at (9.4,1.0)
       {VMX non-root operation};
  \node[b, draw=secondary, fill=white, text=secondary!80!black] at (9.4,0.35)
       {guest kernel --- \textbf{ring 0}};
  \node[b, draw=secondary!50, fill=white, text=secondary!70] at (9.4,-0.3)
       {guest applications --- ring 3};
  \node[lb, text=accent] at (9.4,-0.95)
       {all four rings present; the guest\\cannot tell it is a guest};

  \draw[-{Stealth[length=2.4mm]}, greenhl!70, line width=1.2pt]
        (2.9,0.55) to[out=25,in=155] (6.5,0.55);
  \node[lb, text=greenhl!55!black] at (4.7,1.25)
       {\textbf{VM entry} --- load guest state\\from the VMCS};

  \draw[-{Stealth[length=2.4mm]}, accent, line width=1.2pt]
        (6.5,-0.75) to[out=-155,in=-25] (2.9,-0.75);
  \node[lb, text=accent!70!black] at (4.7,-1.55)
       {\textbf{VM exit} --- save guest state,\\write the exit reason};

  \node[rounded corners=2pt, draw=goldhl, fill=goldhl!10, text=goldhl!55!black,
        font=\scriptsize, align=center, minimum width=3.2cm, inner sep=3pt]
       at (4.7,-2.75)
       {\textbf{VMCS} --- one per virtual CPU\\guest state $\bullet$ host state
        $\bullet$ controls};
  \draw[goldhl!70, line width=0.8pt, dashed] (4.7,-2.2) -- (4.7,-1.95);
\end{tikzpicture}}%
{VT-x in one picture. The guest runs in ring~0 of a second operating mode, so
nothing it executes behaves unexpectedly; the hypervisor regains control
through a hardware transition rather than a trap it had to
arrange.}{vmx-modes}

\begin{conceptbox}[AMD Did the Same Thing with Different Names]
AMD-V (originally \emph{Pacifica}, architecturally \term{SVM}) solves the same
problem with the same structure and an incompatible vocabulary. Guest mode
corresponds to non-root operation; the \term{VMCB} (virtual machine control
block) corresponds to the VMCS; \texttt{VMRUN} corresponds to
\texttt{VMLAUNCH}; \term{nested page tables} correspond to \emph{extended} page
tables.

\smallskip
The differences that matter in practice are small but real: the VMCB is a
plain 4 KB structure a hypervisor reads and writes with ordinary memory
instructions, whereas the VMCS is opaque and must be accessed through
\texttt{VMREAD} and \texttt{VMWRITE}, which cost cycles. AMD's design is
simpler to program; Intel's leaves the processor free to cache VMCS fields in
whatever internal format it likes.

\smallskip
This book uses the Intel names throughout and notes the AMD equivalent where it
differs. Every hypervisor in \chref{platforms} supports both, through an
abstraction layer written once.
\end{conceptbox}

\section{The VMCS, and What a Transition Costs}

A \term{VMCS} --- virtual machine control structure --- exists per virtual CPU
and holds everything needed to switch between the two worlds.

\rowcolors{2}{white}{lightbg}
\begin{longtable}{@{}L{3.6cm}L{11.0cm}@{}}
\toprule
\rowcolor{primary!15}
\textbf{VMCS area} & \textbf{What it holds} \\
\midrule
\endhead
\textbf{Guest-state area} & The guest's registers, control registers, segment
registers, \texttt{RIP} and \texttt{RSP}. Saved automatically on exit, loaded
automatically on entry \\
\textbf{Host-state area} & The hypervisor's corresponding state, loaded on
exit. This is what makes the return to the hypervisor a single hardware
operation rather than a software convention \\
\textbf{VM-execution controls} & \emph{Which events cause an exit.} A bitmap
per instruction class, an exception bitmap, an I/O-port bitmap and an MSR
bitmap. This is the dial the hypervisor turns to trade control against speed \\
\textbf{VM-exit information} & Written by the processor on every exit: the
\term{exit reason}, the faulting address, the instruction length. The
hypervisor reads this first and dispatches on it \\
\bottomrule
\end{longtable}

\begin{conceptbox}[The Execution Controls Are the Real Design Surface]
A novice reading of VT-x is that the processor decides what to trap. It does
not: the \emph{hypervisor} decides, by setting bits in the execution controls,
and the engineering is in choosing well.

\smallskip
Exit on every \texttt{CPUID}? You must, because the guest must not see the
host's true feature set. Exit on every \texttt{RDTSC}? Only if you care about
hiding the passage of time, and most hypervisors do not, because the exit cost
would be ruinous. Exit on every write to \texttt{CR3}? With shadow page tables
you must; with EPT you need not, and that single bit is most of the next
section.

\smallskip
Every bit set buys correctness or isolation and costs an exit. Hypervisor
performance work is, to a first approximation, the work of clearing bits that
turn out not to be needed.
\end{conceptbox}

\begin{alertbox}[A VM exit is a world switch, not a function call]
On exit the processor must save the guest's entire architectural state, load
the hypervisor's, and flush or tag anything that was speculating on the guest's
address space. On the first VT-x parts this cost on the order of 1{,}500 cycles
in each direction; a decade of optimisation brought it to a few hundred, and it
has not fallen much since because the work is irreducible.

\smallskip
That number is why \chref{architectures} could report that binary translation
beat early hardware assistance: a translated call into the monitor cost tens of
cycles and did not disturb the pipeline. Hardware assistance won in the end not
by making exits cheap but by \emph{removing the reasons to exit} --- which is
what extended page tables and the MSR bitmap are for.
\end{alertbox}

\section{Two-Dimensional Paging}

Virtualizing the processor was half the problem. The other half is memory, and
it is harder.

A guest kernel maintains page tables mapping \emph{guest virtual} addresses to
what it believes are physical addresses. They are not physical: they are
\term{guest physical} addresses, which the hypervisor must in turn map to
\term{host physical} addresses. Two translations, and the hardware originally
performed only one.

\subsection{The first answer: shadow page tables}

The hypervisor builds its own page tables mapping guest virtual directly to
host physical --- \term{shadow page tables} --- and points the real \texttt{CR3}
at those. The processor walks them normally, so a TLB miss costs exactly what
it costs natively.

The price is keeping them consistent. The hypervisor must trap every guest
write to its own page tables, which means write-protecting them and taking a VM
exit on each modification. A process fork, which rewrites a great many page
table entries, becomes thousands of exits.

\subsection{The second answer: let the hardware do both}

\term{Extended page tables} (Intel) and \term{nested page tables} (AMD) add a
second, hardware-walked translation. The guest's own page tables stay in
charge of guest-virtual to guest-physical, and the guest may edit them freely
with no exits at all. A second set of tables, owned by the hypervisor, maps
guest-physical to host-physical, and the processor walks both.

\begin{worked}[Counting a Nested Page Walk]
How much does the hardware actually do on a TLB miss under EPT? Count memory
accesses for 4-level paging on x86-64.

\smallskip
\textbf{Native, for comparison.} Four levels --- PML4, PDPT, PD, PT --- so a
TLB miss costs \textbf{4 memory accesses}.

\smallskip
\textbf{Under EPT.} Every address the walker touches is a \emph{guest
physical} address, and each one must itself be translated by a 4-level EPT
walk before it can be used.

\begin{tight}
  \item Guest \texttt{CR3} holds the GPA of the guest PML4. Translate it:
        \textbf{4} EPT accesses. Then read the guest PML4 entry: \textbf{1}.
  \item That entry holds the GPA of the PDPT. Translate it: \textbf{4}. Read
        the PDPT entry: \textbf{1}.
  \item GPA of the PD. Translate: \textbf{4}. Read the PD entry: \textbf{1}.
  \item GPA of the PT. Translate: \textbf{4}. Read the PT entry: \textbf{1}.
  \item The guest PTE yields the final GPA of the data page. Translate it:
        \textbf{4}.
\end{tight}

\smallskip
\textbf{The total.}
\[ \underbrace{4 \times (4 + 1)}_{\text{four guest levels}} \;+\;
   \underbrace{4}_{\text{final translation}} \;=\; \mathbf{24}
   \text{ memory accesses} \]

\smallskip
\textbf{So EPT is six times worse than native on a TLB miss}, and worse than
shadow paging, which costs 4. If that were the whole story, shadow page tables
would have won.

\smallskip
\textbf{Why EPT won anyway.} Count the \emph{other} operation. Under shadow
paging every guest page-table write is a VM exit at several hundred cycles;
under EPT it is an ordinary memory write at a few cycles. A workload that
creates processes, maps files or frees memory does this constantly.

\smallskip
Measured on a fork-heavy build: shadow paging took tens of thousands of exits
per second; EPT took almost none. The 24-access walk is also mitigated in
hardware by caching the intermediate EPT translations --- a page-walk cache ---
so the full 24 is the worst case rather than the common one.

\smallskip
\textbf{The lesson.} The right structure is not the one with the better
best case; it is the one whose expensive case is rare. Shadow paging is
cheaper per TLB miss and catastrophic per page-table edit. EPT is the reverse,
and page-table edits are far more common than the arithmetic suggests.
\end{worked}

\dsfig{%
\begin{tikzpicture}[font=\scriptsize,
  lvl/.style={rounded corners=2pt, line width=0.7pt, align=center,
              minimum width=1.5cm, minimum height=0.5cm, inner sep=2pt},
  g/.style={lvl, draw=secondary, fill=secondary!10, text=secondary!80!black},
  e/.style={lvl, draw=purplehl, fill=purplehl!10, text=purplehl!70!black,
            minimum width=0.62cm, minimum height=0.42cm, font=\tiny},
  ar/.style={-{Stealth[length=1.6mm]}, line width=0.7pt}
]
  \node[font=\scriptsize\bfseries, text=secondary!80!black, anchor=east]
       at (-0.95,3.3) {guest walk};
  \node[font=\scriptsize\bfseries, text=purplehl!70!black, anchor=east]
       at (-0.95,2.45) {EPT walk for each};

  \foreach \i/\n in {0/{guest\\PML4}, 1/{guest\\PDPT}, 2/{guest\\PD},
                     3/{guest\\PT}, 4/{data\\page}}{
    \pgfmathsetmacro{\gx}{\i*3.0}
    \node[g] at (\gx,3.3) {\n};
    \foreach \j in {0,1,2,3}{
      \pgfmathsetmacro{\ex}{\gx-0.93+\j*0.62}
      \node[e] at (\ex,2.45) {};}
    \draw[ar, purplehl!60] (\gx,2.72) -- (\gx,3.05);}

  \foreach \i in {0,1,2,3}{
    \pgfmathsetmacro{\ax}{\i*3.0+0.78}
    \pgfmathsetmacro{\bx}{\i*3.0+2.22}
    \draw[ar, secondary!70] (\ax,3.3) -- (\bx,3.3);}

  \node[font=\scriptsize\itshape, text=accent, anchor=west] at (-1.4,1.75)
       {each purple block is 4 memory accesses $\bullet$
        4 guest reads $+$ 5 EPT walks of 4 $=$ \textbf{24} accesses,
        against 4 natively};

  % --- the comparison bar ---
  \node[font=\scriptsize\bfseries, text=neutral!40!black, anchor=east]
       at (-0.95,0.85) {TLB miss};
  \node[font=\scriptsize\bfseries, text=neutral!40!black, anchor=east]
       at (-0.95,0.25) {page-table write};
  \foreach \y/\lab/\sh/\ep in {0.85/{}/{4 accesses}/{24 accesses},
                               0.25/{}/{VM exit, ~500 cycles}/{a few cycles}}{
    \node[rounded corners=2pt, draw=goldhl, fill=goldhl!8, text=goldhl!55!black,
          minimum width=4.4cm, minimum height=0.45cm, font=\scriptsize]
         at (1.6,\y) {shadow: \sh};
    \node[rounded corners=2pt, draw=purplehl, fill=purplehl!8,
          text=purplehl!70!black, minimum width=4.4cm, minimum height=0.45cm,
          font=\scriptsize] at (6.6,\y) {EPT: \ep};}
  \node[font=\scriptsize\itshape, text=neutral, anchor=west, align=left]
       at (9.2,0.55) {EPT loses the rare case\\and wins the common one};
\end{tikzpicture}}%
{A nested page walk, and the trade it represents. Shadow paging is cheaper on a
TLB miss and ruinous on every guest page-table edit; extended page tables
invert both, and edits are the more frequent event.}{nested-paging}

\begin{conceptbox}[VPID: Not Flushing the TLB Every Time]
One more cost was hiding in the transitions. Early VT-x flushed the entire TLB
on every VM entry and exit, because the hypervisor and the guest use different
address spaces and the processor had no way to tell their entries apart. A busy
guest therefore began every re-entry with a cold TLB.

\smallskip
\term{VPID} --- virtual processor identifier, AMD's equivalent is the ASID ---
tags each TLB entry with the virtual CPU it belongs to. Entries survive the
transition and are simply ignored while another context runs. It is a small
feature with a large effect, and it is the clearest illustration of the
pattern in this chapter: the exits were not made cheaper, the work done around
them was removed.
\end{conceptbox}

\section{The IOMMU}

The processor is virtualized and memory is virtualized. Devices are not, and
devices are dangerous, because a device performing DMA writes to \emph{physical}
addresses and knows nothing about any of this.

Hand a network card directly to a guest --- which \chref{paravirt} will want to
do for speed --- and the guest's driver programs it with addresses from its own
guest-physical space. The card would write wherever those happen to land in
real memory: another guest, or the hypervisor. A single passthrough device
without protection destroys the resource-control property for the entire host.

\begin{definitionbox}[IOMMU, DMA Remapping, Interrupt Remapping]
An \term{IOMMU} --- Intel calls the feature \term{VT-d}, AMD calls it AMD-Vi ---
is a memory management unit that sits between devices and memory, as the MMU
sits between the processor and memory.

\smallskip
\term{DMA remapping}: every address a device emits is translated through
per-device page tables that the hypervisor controls. A device assigned to a
guest can reach that guest's memory and nothing else. An address outside the
mapping faults rather than being written.

\smallskip
\term{Interrupt remapping}: a device cannot raise an arbitrary interrupt vector
on an arbitrary core. Interrupts are translated too, and delivered only where
the hypervisor allows.

\smallskip
Without an IOMMU, device passthrough is not merely risky; it is equivalent to
giving the guest write access to all of physical memory.
\end{definitionbox}

\begin{alertbox}[Nested virtualization, and why your lab needs it]
Running a hypervisor \emph{inside} a guest --- nested virtualization --- means
the processor must present VMX itself as a virtualizable feature. The outer
hypervisor intercepts the inner one's \texttt{VMLAUNCH}, merges the inner
VMCS with its own, and runs the result. It works, and it is how the labs in
this book run a Type-1 hypervisor on a machine that is already a guest.

\smallskip
It is also slow: every exit the inner guest takes may become several exits in
the outer one. Expect a nested guest to be noticeably slower than a
single-level one, and do not use nesting to measure anything.

\smallskip
On a laptop it must usually be enabled explicitly. Appendix~A gives the
commands; the lab below checks it first, because every later chapter's lab
assumes it.
\end{alertbox}

\begin{pitfall}
\begin{tight}
  \item \textbf{Believing VT-x made virtualization fast.} It made it
        \emph{correct} --- the architecture now satisfies Popek--Goldberg. The
        first generation was often slower than binary translation. Speed came
        from EPT, VPID and the bitmaps, over roughly a decade.
  \item \textbf{Assuming EPT is faster at everything.} It is six times slower
        than native on a TLB miss and six times slower than shadow paging. It
        wins on the workload, not on the microbenchmark, and a workload with a
        huge working set and few page-table edits is exactly where shadow
        paging would still win.
  \item \textbf{Confusing the MMU with the IOMMU.} The MMU protects memory from
        the processor; the IOMMU protects memory from \emph{devices}. A host
        with VT-x but no VT-d can virtualize fine and cannot safely pass a
        device through.
  \item \textbf{Leaving virtualization disabled in firmware and blaming the
        software.} A large share of ``KVM will not start'' is a cleared BIOS
        bit. Check \texttt{/proc/cpuinfo} first --- the lab's step 1 is there
        for this reason.
  \item \textbf{Benchmarking inside a nested guest.} Every measurement is
        inflated by the outer hypervisor's exits, and by a variable amount.
        Nesting is for learning, not for numbers.
\end{tight}
\end{pitfall}

\begin{lab}[Watching the Hardware Work]
Counts real VM exits on a real guest and shows what changes them. Needs KVM and
a Linux host; \texttt{perf} comes from \texttt{linux-tools}. Expect thirty
minutes.

\begin{lstlisting}[language=bash]
#!/usr/bin/env bash
# ch04_exits.sh -- see VMX, EPT and the exit rate with your own eyes.
set -euo pipefail

# --- 1. Does this processor have the features this chapter describes? ---
grep -o -m1 -E '\<(vmx|svm)\>' /proc/cpuinfo    # VT-x or AMD-V at all
grep -o -m1 -E '\<ept\>|\<npt\>' /proc/cpuinfo  # two-dimensional paging
grep -o -m1 '\<vpid\>' /proc/cpuinfo            # tagged TLB entries
# VT-d / AMD-Vi shows up as IOMMU groups once enabled in firmware:
ls /sys/kernel/iommu_groups 2>/dev/null | wc -l

# --- 2. Is nested virtualization on? Every later lab wants it. ---------
cat /sys/module/kvm_intel/parameters/nested 2>/dev/null \
  || cat /sys/module/kvm_amd/parameters/nested 2>/dev/null
# If that prints N or 0, see Appendix A for how to turn it on.

# --- 3. Start the guest from Chapter 1 and find its thread. ------------
sudo virsh start vss1 || true
sleep 25
PID=$(pgrep -f 'qemu.*vss1' | head -1)
echo "guest qemu pid: $PID"

# --- 4. Count exits by reason, for ten seconds of an idle guest. -------
# Every line is a different entry in the VM-execution controls.
sudo perf stat -e 'kvm:kvm_exit' -p "$PID" -- sleep 10
sudo perf record -e 'kvm:kvm_exit' -p "$PID" -- sleep 10 2>/dev/null
sudo perf script \
  | awk '{for(i=1;i<=NF;i++) if($i ~ /reason/) print $(i+1)}' \
  | sort | uniq -c | sort -rn | head -12

# --- 5. Now give the guest work that forces exits, and compare. --------
# Network I/O through an emulated device is an exit per register access.
in_guest () { ssh -o StrictHostKeyChecking=no "ubuntu@$1" "$2"; }
IP=$(sudo virsh domifaddr vss1 | awk '/ipv4/{split($4,a,"/"); print a[1]}')
in_guest "$IP" "ping -f -c 2000 10.0.2.2 >/dev/null 2>&1" &
sudo perf stat -e 'kvm:kvm_exit' -p "$PID" -- sleep 10
wait

# --- 6. Watch EPT doing its job: no exit for a guest page-table edit. --
# Fork a few thousand processes; under shadow paging this storms the host.
in_guest "$IP" "for i in \$(seq 2000); do /bin/true; done"
sudo perf stat -e 'kvm:kvm_exit' -p "$PID" -- sleep 5

# --- 7. Tear down ------------------------------------------------------
sudo virsh destroy vss1 || true
\end{lstlisting}

\textbf{What to notice.} The idle guest still exits, a few thousand times a
second, almost all of it timer and interrupt work. The flood ping multiplies
that by a large factor, and the two thousand forks barely move it --- which is
the worked example's claim, measured. Write the three numbers down.
\end{lab}

\begin{checkpoint}
\begin{tightnum}
  \item Why did Intel and AMD choose to add a new processor mode rather than
        simply making the seventeen offending instructions privileged?
  \item Under 5-level paging, a guest page walk has five levels and so does
        the EPT walk. How many memory accesses does a nested walk cost, and
        what is the general formula for an $n$-level scheme?
  \item A host has VT-x and EPT but the IOMMU is disabled in firmware. State
        exactly which capability is lost, and why the hypervisor must refuse
        it rather than merely warn.
\end{tightnum}
\end{checkpoint}

\begin{chaptersummary}
\begin{tight}
  \item VT-x and AMD-V add a second operating mode rather than changing any
        instruction, so the guest kernel runs in ring~0 and backward
        compatibility is preserved. x86 now satisfies the Popek--Goldberg
        condition.
  \item The VMCS holds guest state, host state, the controls that decide what
        causes an exit, and the exit information. The controls are the design
        surface: every bit set buys control and costs exits.
  \item A VM exit is a world switch --- around 1{,}500 cycles originally, a few
        hundred now. Hardware assistance won by removing reasons to exit, not
        by making exits cheap.
  \item Shadow page tables cost 4 accesses per TLB miss and a VM exit per guest
        page-table write. Extended page tables cost 24 accesses per miss and
        nothing per write. Writes are commoner, so EPT won.
  \item VPID tags TLB entries so they survive a transition.
  \item An IOMMU translates and confines device DMA. Without one, passthrough
        gives a guest write access to all physical memory.
  \item Nested virtualization works and is slow. Use it to learn, never to
        measure.
\end{tight}
\end{chaptersummary}

\begin{reviewq}
  \item \textbf{(MCQ)} Under VT-x a guest kernel runs in: (a)~ring~1 of root
        operation; (b)~ring~0 of non-root operation; (c)~ring~3 of non-root
        operation; (d)~ring~0 of root operation.
  \item \textbf{(MCQ)} Extended page tables principally reduce: (a)~the cost of
        a TLB miss; (b)~the number of VM exits caused by guest memory
        management; (c)~the size of the VMCS; (d)~interrupt latency.
  \item \textbf{(Short)} Explain what the VM-execution controls are for, and
        give one example of a bit you would set and one you would clear, with
        the reason for each.
  \item \textbf{(Short)} Derive the 24-access figure for a nested page walk,
        showing where each group of accesses comes from.
  \item \textbf{(Short)} State what an IOMMU protects against, and describe the
        concrete failure that occurs without one.
  \item \textbf{(Applied)} A guest is slow and \texttt{perf} shows 180{,}000 VM
        exits per second, 80\% of them with an I/O-instruction exit reason.
        Explain what the guest is most likely doing, name the change you would
        make, and predict what the exit count would become and why.
  \item \textbf{(Applied)} You are asked whether a six-year-old server is
        suitable as a teaching host for this course. List the processor and
        firmware features you would check, in order, and say which lab in this
        book each one is needed for.
\end{reviewq}
